\section{Cofinal sequences of regular cardinals}

\import{function.tex}
\import{cardinal/infinite.tex}
\import{cardinal/cofinality.tex}
\import{cardinal/normal-map.tex}
\import{cardinal/normal-sequence.tex}
\import{cardinal/successor.tex}
\import{cardinal/regularity.tex}

A cofinal regular-cardinal sequence is required only to be strictly
increasing.  It need not be continuous: taking successor cardinals generally
destroys continuity at limit indices.

\begin{definition}\label{cofinal_regular_cardinal_sequence}
    $\cofinalRegularCardinalSequence{r}{\lambda}{\kappa}$ iff
        $\lambda$ and $\kappa$ are ordinals and
        $\cofinalMap{r}{\lambda}{\kappa}$ and
        $\normalMapIncreasing{r}{\lambda}$ and
        for every $\xi\in\lambda$ we have
            $r(\xi)$ is a regular cardinal.
\end{definition}

\begin{proposition}\label{cofinal_regular_cardinal_sequence_type}
    Suppose $\cofinalRegularCardinalSequence{r}{\lambda}{\kappa}$.
    Then $\lambda$ and $\kappa$ are ordinals and
        $r\in\funs{\lambda}{\kappa}$.
\end{proposition}
\begin{proof}
    $\lambda$ and $\kappa$ are ordinals and
        $\cofinalMap{r}{\lambda}{\kappa}$ by
        \obv{cofinal_regular_cardinal_sequence}.
    Hence $r\in\funs{\lambda}{\kappa}$ by \obv{cofinal_map}.
\end{proof}

\begin{proposition}\label{cofinal_regular_cardinal_sequence_is_cofinal}
    Suppose $\cofinalRegularCardinalSequence{r}{\lambda}{\kappa}$.
    Then $\cofinalMap{r}{\lambda}{\kappa}$.
\end{proposition}
\begin{proof}
    Follows by definition.
\end{proof}

\begin{proposition}\label{cofinal_regular_cardinal_sequence_increasing}
    Suppose $\cofinalRegularCardinalSequence{r}{\lambda}{\kappa}$.
    Let $\eta\in\lambda$ and $\xi\in\eta$.
    Then $r(\xi)\in r(\eta)$.
\end{proposition}
\begin{proof}
    $\normalMapIncreasing{r}{\lambda}$ by
        \obv{cofinal_regular_cardinal_sequence}.
    Follows by \obv{normal_map_increasing}.
\end{proof}

\begin{proposition}\label{cofinal_regular_cardinal_sequence_value}
    Suppose $\cofinalRegularCardinalSequence{r}{\lambda}{\kappa}$.
    Let $\xi\in\lambda$.
    Then $r(\xi)$ is a regular cardinal.
\end{proposition}
\begin{proof}
    Follows by definition.
\end{proof}

\begin{proposition}\label{cofinal_regular_cardinal_sequence_value_is_infinite_cardinal}
    Suppose $\cofinalRegularCardinalSequence{r}{\lambda}{\kappa}$.
    Let $\xi\in\lambda$.
    Then $r(\xi)$ is an infinite cardinal.
\end{proposition}
\begin{proof}
    $r(\xi)$ is a regular cardinal by
        \obv{cofinal_regular_cardinal_sequence_value}.
    Follows by \obv{regular_cardinal}.
\end{proof}

\begin{proposition}\label{cofinal_regular_cardinal_sequence_value_is_cardinal}
    Suppose $\cofinalRegularCardinalSequence{r}{\lambda}{\kappa}$.
    Let $\xi\in\lambda$.
    Then $r(\xi)$ is a cardinal.
\end{proposition}
\begin{proof}
    $r(\xi)$ is an infinite cardinal by
        \obv{cofinal_regular_cardinal_sequence_value_is_infinite_cardinal}.
    Follows by definition.
\end{proof}

\begin{proposition}\label{cofinal_regular_cardinal_sequence_value_is_ordinal}
    Suppose $\cofinalRegularCardinalSequence{r}{\lambda}{\kappa}$.
    Let $\xi\in\lambda$.
    Then $r(\xi)$ is an ordinal.
\end{proposition}
\begin{proof}
    $r(\xi)$ is a cardinal by
        \obv{cofinal_regular_cardinal_sequence_value_is_cardinal}.
    Follows by \auto{cardinal,initial_ordinal}.
\end{proof}

\begin{theorem}\label{cofinal_regular_sequence_reindex}
    Suppose $\cofinalRegularCardinalSequence{r}{\lambda}{\kappa}$.
    Let $\mu$ be an ordinal.
    Suppose $\cofinalMap{q}{\mu}{\lambda}$ and
        $\normalMapIncreasing{q}{\mu}$.
    Then
        $\cofinalRegularCardinalSequence{r\circ q}{\mu}{\kappa}$.
\end{theorem}
\begin{proof}
    $\lambda$ and $\kappa$ are ordinals and
        $r\in\funs{\lambda}{\kappa}$ by
        \obv{cofinal_regular_cardinal_sequence_type}.
    $q\in\funs{\mu}{\lambda}$ by \obv{cofinal_map}.
    Hence $q$ and $r$ are functions and
        $\dom{q}=\mu$ and $\dom{r}=\lambda$ by
        \obv{funs_is_function,funs}.
    $\ran{q}\subseteq\lambda$ by \obv{funs_ran}.
    Hence $\ran{q}\subseteq\dom{r}$.
    Hence $r$ is composable with $q$ by definition.
    $r\circ q\in\funs{\mu}{\kappa}$ by \obv{funs_circ}.
    Show that $\normalMapIncreasing{r\circ q}{\mu}$.
    \begin{subproof}
        Show that for all $i,j$ such that $i\in j\in\mu$ we have
            $(r\circ q)(i)\in(r\circ q)(j)$.
        \begin{subproof}
            Fix $i,j$.
            Assume $i\in j\in\mu$.
            Hence $i,j\in\mu$ by \obv{ordinal_transitivity}.
            Hence $i,j\in\dom{q}$.
            Hence $q(i)\in q(j)$ by \obv{normal_map_increasing}.
            $q(j)\in\lambda$ by \obv{funs_type_apply}.
            Hence $r(q(i))\in r(q(j))$ by
                \obv{cofinal_regular_cardinal_sequence_increasing}.
            $(r\circ q)(i)=r(q(i))$ by \obv{circ_apply}.
            $(r\circ q)(j)=r(q(j))$ by \obv{circ_apply}.
        \end{subproof}
        Follows by \obv{normal_map_increasing}.
    \end{subproof}
    Show that for every $i\in\mu$ we have
        $(r\circ q)(i)$ is a regular cardinal.
    \begin{subproof}
        Fix $i\in\mu$.
        Hence $i\in\dom{q}$.
        $q(i)\in\lambda$ by \obv{funs_type_apply}.
        Hence $r(q(i))$ is a regular cardinal by
            \obv{cofinal_regular_cardinal_sequence_value}.
        $(r\circ q)(i)=r(q(i))$ by \obv{circ_apply}.
    \end{subproof}
    Show that $\cofinalMap{r\circ q}{\mu}{\kappa}$.
    \begin{subproof}
        $\cofinalIn{\ran{r}}{\kappa}$ and
            $\cofinalIn{\ran{q}}{\lambda}$ by
            \auto{cofinal_regular_cardinal_sequence_is_cofinal,cofinal_map}.
        Show that for every $\alpha\in\kappa$ there exists
            $\delta\in\ran{r\circ q}$ such that $\alpha\subseteq\delta$.
        \begin{subproof}
            Fix $\alpha\in\kappa$.
            Take $\gamma\in\ran{r}$ such that $\alpha\subseteq\gamma$ by
                \obv{cofinal_in}.
            Take $\xi\in\dom{r}$ such that $r(\xi)=\gamma$ by
                \obv{fun_ran_iff}.
            Hence $\xi\in\lambda$.
            Take $\zeta\in\ran{q}$ such that $\xi\subseteq\zeta$ by
                \obv{cofinal_in}.
            Take $i\in\dom{q}$ such that $q(i)=\zeta$ by
                \obv{fun_ran_iff}.
            Hence $i\in\mu$.
            $q(i)\in\lambda$ by \obv{funs_type_apply}.
            Hence $\xi\subseteq q(i)$.
            $\xi$ is an ordinal by \obv{elem_of_ordinal_is_ordinal}.
            $q(i)$ is an ordinal by \obv{elem_of_ordinal_is_ordinal}.
            Hence $\xi\in q(i)$ or $\xi=q(i)$ by
                \obv{ordinal_subseteq_implies_elem_or_eq}.
            Show that $r(\xi)\subseteq r(q(i))$.
            \begin{subproof}
                \begin{byCase}
                    \caseOf{$\xi=q(i)$.}
                        Follows by \obv{subseteq_refl}.
                    \caseOf{$\xi\neq q(i)$.}
                        Hence $\xi\in q(i)$.
                        Hence $r(\xi)\in r(q(i))$ by
                            \obv{cofinal_regular_cardinal_sequence_increasing}.
                        $r(\xi)$ and $r(q(i))$ are ordinals by
                            \auto{cofinal_regular_cardinal_sequence_value,
                                regular_cardinal,cardinal,initial_ordinal}.
                        Follows by \obv{ordinal_elem_implies_subseteq}.
                \end{byCase}
            \end{subproof}
            Hence $\alpha\subseteq r(q(i))$ by
                \obv{subseteq_transitive}.
            $(r\circ q)(i)=r(q(i))$ by \obv{circ_apply}.
            $(r\circ q)(i)\in\ran{r\circ q}$ by
                \auto{fun_ran_iff,funs_is_function,funs}.
        \end{subproof}
        $\ran{r\circ q}\subseteq\kappa$ by \obv{funs_ran}.
        Hence $\cofinalIn{\ran{r\circ q}}{\kappa}$ by
            \obv{cofinal_in}.
        Follows by \obv{cofinal_map}.
    \end{subproof}
    Follows by \obv{cofinal_regular_cardinal_sequence}.
\end{proof}

\begin{corollary}\label{cofinal_regular_sequence_normal_reindex}
    Suppose $\cofinalRegularCardinalSequence{r}{\lambda}{\kappa}$ and
        $\normalCofinalMap{q}{\mu}{\lambda}$.
    Then
        $\cofinalRegularCardinalSequence{r\circ q}{\mu}{\kappa}$.
\end{corollary}
\begin{proof}
    $\mu$ is an ordinal by \auto{normal_cofinal_map_type}.
    $\cofinalMap{q}{\mu}{\lambda}$ by
        \obv{normal_cofinal_map_is_cofinal}.
    $\normalMap{q}{\mu}{\lambda}$ by
        \obv{normal_cofinal_map_is_normal}.
    Hence $\normalMapIncreasing{q}{\mu}$ by \obv{normal_map}.
    Follows by \obv{cofinal_regular_sequence_reindex}.
\end{proof}

\subsection{Successor-cardinal regularization}

\begin{definition}\label{successor_cardinal_map}
    $\successorCardinalMap{c}{\lambda}
        =\{(\xi,\cardSucc{c(\xi)})\mid\xi\in\lambda\}$.
\end{definition}

\begin{proposition}\label{successor_cardinal_map_iff}
    $(\xi,\rho)\in\successorCardinalMap{c}{\lambda}$ iff
        $\xi\in\lambda$ and $\rho=\cardSucc{c(\xi)}$.
\end{proposition}
\begin{proof}
    Follows by \auto{successor_cardinal_map,pair_eq_iff}.
\end{proof}

\begin{theorem}\label{successor_cardinal_map_type}
    Suppose for every $\xi\in\lambda$ we have
        $\cardSucc{c(\xi)}\in\kappa$.
    Then $\successorCardinalMap{c}{\lambda}
        \in\funs{\lambda}{\kappa}$.
\end{theorem}
\begin{proof}
    Let $r=\successorCardinalMap{c}{\lambda}$.
    $r$ is a relation by
        \auto{successor_cardinal_map,relation}.
    $r$ is right-unique by
        \auto{successor_cardinal_map,rightunique,pair_eq_iff}.
    $\dom{r}=\lambda$ by
        \auto{setext,dom_iff,successor_cardinal_map,pair_eq_iff}.
    For every $\xi\in\dom{r}$ we have
        $r(\xi)=\cardSucc{c(\xi)}$ by
        \auto{successor_cardinal_map,function_apply_intro}.
    Hence $r$ is a function to $\kappa$ by \auto.
    Follows by \auto{funs_intro}.
\end{proof}

\begin{proposition}\label{successor_cardinal_map_apply}
    Let $\xi\in\lambda$.
    Then $(\successorCardinalMap{c}{\lambda})(\xi)
        =\cardSucc{c(\xi)}$.
\end{proposition}
\begin{proof}
    $\successorCardinalMap{c}{\lambda}$ is a function by
        \auto{successor_cardinal_map,relation,rightunique,pair_eq_iff}.
    $(\xi,\cardSucc{c(\xi)})
        \in\successorCardinalMap{c}{\lambda}$ by
        \obv{successor_cardinal_map}.
    Follows by \obv{function_apply_intro}.
\end{proof}

\begin{theorem}\label{cofinal_regular_sequence_above_index_exists}
    Let $\kappa$ be a singular cardinal.
    Then there exists $r$ such that
        $\cofinalRegularCardinalSequence
            {r}{\cofinality{\kappa}}{\kappa}$ and
        $\cofinality{\kappa}\in r(\emptyset)$.
\end{theorem}
\begin{proof}
    Hence $\kappa$ is a limit ordinal by
        \obv{infinite_cardinal_is_limit_ordinal}.
    Take $d$ such that
        $\cofinalMap{d}{\cofinality{\kappa}}{\kappa}$ by
        \obv{cofinality_map_exists}.
    $\normalCofinalCardinalSequence
            {\cardinalSequence{\kappa}{d}}{\cofinality{\kappa}}{\kappa}$ and
        $\infiniteCardinalSequenceValues
            {\cardinalSequence{\kappa}{d}}{\cofinality{\kappa}}$ by
        \auto{cardinal_sequence_is_normal_cofinal}.
    $\cofinality{\kappa}
        \in(\cardinalSequence{\kappa}{d})(\emptyset)$ by
        \obv{cardinal_sequence_contains_cofinality_at_zero}.
    Let $\lambda=\cofinality{\kappa}$.
    $\cofinalMap{d}{\lambda}{\kappa}$.
    $\lambda$ is a regular cardinal by
        \obv{cofinality_of_limit_is_regular}.
    Hence $\lambda$ is a limit ordinal by
        \obv{infinite_cardinal_is_limit_ordinal}.
    Hence $\lambda$ is an ordinal and $\emptyset\in\lambda$ by
        \obv{limit_ordinal}.
    Let $c=\cardinalSequence{\kappa}{d}$.
    $\normalCofinalCardinalSequence{c}{\lambda}{\kappa}$ and
        $\infiniteCardinalSequenceValues{c}{\lambda}$ and
        $\lambda\in c(\emptyset)$.
    Hence for every $\xi\in\lambda$ we have
        $c(\xi)$ is an infinite cardinal by
        \obv{infinite_cardinal_sequence_value}.
    Hence $\normalCofinalMap{c}{\lambda}{\kappa}$ by
        \obv{normal_cofinal_cardinal_sequence_is_normal_cofinal}.
    Hence $\cofinalMap{c}{\lambda}{\kappa}$ and
        $c\in\funs{\lambda}{\kappa}$ by
        \obv{normal_cofinal_map_is_cofinal,normal_cofinal_map_type}.
    Hence $c$ is a function and $\dom{c}=\lambda$ by
        \obv{funs_is_function,funs}.
    Show that for every $\xi\in\lambda$ we have
        $\cardSucc{c(\xi)}\in\kappa$.
    \begin{subproof}
        Fix $\xi\in\lambda$.
        $c(\xi)\in\kappa$ by \obv{funs_type_apply}.
        Follows by
            \obv{infinite_cardinal_successor_belongs_to_singular_cardinal}.
    \end{subproof}
    Hence $\successorCardinalMap{c}{\lambda}
        \in\funs{\lambda}{\kappa}$ by
        \obv{successor_cardinal_map_type}.
    Let $r=\successorCardinalMap{c}{\lambda}$.
    $r\in\funs{\lambda}{\kappa}$.
    Hence $r$ is a function by \obv{funs_is_function}.
    $\dom{r}=\lambda$ by \obv{funs}.
    Show that for every $\xi\in\lambda$ we have
        $r(\xi)$ is a regular cardinal.
    \begin{subproof}
        Fix $\xi\in\lambda$.
        $c(\xi)$ is an infinite cardinal.
        Hence $\cardSucc{c(\xi)}$ is a regular cardinal by
            \obv{infinite_cardinal_successor_is_regular}.
        $r(\xi)=\cardSucc{c(\xi)}$ by
            \auto{successor_cardinal_map_apply}.
    \end{subproof}
    Show that $\normalMapIncreasing{r}{\lambda}$.
    \begin{subproof}
        Show that for all $\xi,\eta$ such that
            $\xi\in\eta\in\lambda$ we have $r(\xi)\in r(\eta)$.
        \begin{subproof}
            Fix $\xi,\eta$.
            Assume $\xi\in\eta\in\lambda$.
            Hence $\xi,\eta\in\lambda$ by
                \obv{ordinal_transitivity}.
            $c(\xi)$ and $c(\eta)$ are infinite cardinals.
            $c(\xi)\in c(\eta)$ by
                \obv{normal_cofinal_cardinal_sequence_increasing}.
            Hence $\cardLt{c(\xi)}{c(\eta)}$ by
                \auto{card_of_cardinal,card_elem_iff_cardinal_lt}.
            Hence $\cardSucc{c(\xi)}\subseteq c(\eta)$ by
                \obv{cardinal_successor_minimal}.
            Hence $\cardLe{\cardSucc{c(\xi)}}{c(\eta)}$ by
                \obv{subseteq_cardinal_le}.
            $\cardLt{c(\eta)}{\cardSucc{c(\eta)}}$ by
                \obv{cardinal_lt_successor}.
            Hence $\cardLt{\cardSucc{c(\xi)}}
                {\cardSucc{c(\eta)}}$ by
                \obv{cardinal_le_lt_transitive}.
            $\cardSucc{c(\xi)}$ and $\cardSucc{c(\eta)}$ are cardinals by
                \obv{cardinal_successor_is_cardinal}.
            Hence $\cardSucc{c(\xi)}\in\cardSucc{c(\eta)}$ by
                \auto{card_of_cardinal,card_elem_iff_cardinal_lt}.
            $r(\xi)=\cardSucc{c(\xi)}$ and
                $r(\eta)=\cardSucc{c(\eta)}$ by
                \auto{successor_cardinal_map_apply}.
        \end{subproof}
        Follows by \obv{normal_map_increasing}.
    \end{subproof}
    Show that $\cofinalMap{r}{\lambda}{\kappa}$.
    \begin{subproof}
        $\cofinalIn{\ran{c}}{\kappa}$ by \obv{cofinal_map}.
        Show that for every $\alpha\in\kappa$ there exists
            $\rho\in\ran{r}$ such that $\alpha\subseteq\rho$.
        \begin{subproof}
            Fix $\alpha\in\kappa$.
            Take $\gamma\in\ran{c}$ such that $\alpha\subseteq\gamma$ by
                \obv{cofinal_in}.
            Take $\xi\in\dom{c}$ such that $c(\xi)=\gamma$ by
                \obv{fun_ran_iff}.
            Hence $\xi\in\lambda$.
            $c(\xi)$ is an infinite cardinal.
            Hence $c(\xi)$ is an ordinal by
                \auto{cardinal,initial_ordinal}.
            $\cardSucc{c(\xi)}$ is an ordinal by
                \auto{cardinal_successor_is_cardinal,cardinal,initial_ordinal}.
            $c(\xi)\in\cardSucc{c(\xi)}$ by
                \obv{cardinal_belongs_to_successor}.
            Hence $c(\xi)\subseteq\cardSucc{c(\xi)}$ by
                \obv{ordinal_elem_implies_subseteq}.
            Hence $\alpha\subseteq\cardSucc{c(\xi)}$ by
                \obv{subseteq_transitive}.
            $r(\xi)=\cardSucc{c(\xi)}$ by
                \auto{successor_cardinal_map_apply}.
            $r(\xi)\in\ran{r}$ by \auto{fun_ran_iff}.
        \end{subproof}
        $\ran{r}\subseteq\kappa$ by \obv{funs_ran}.
        Hence $\cofinalIn{\ran{r}}{\kappa}$ by \obv{cofinal_in}.
        Follows by \obv{cofinal_map}.
    \end{subproof}
    Hence $\cofinalRegularCardinalSequence{r}{\lambda}{\kappa}$ by
        \obv{cofinal_regular_cardinal_sequence}.
    $c(\emptyset)$ is an infinite cardinal.
    Hence $c(\emptyset)$ is an ordinal by
        \auto{cardinal,initial_ordinal}.
    $\cardSucc{c(\emptyset)}$ is an ordinal by
        \auto{cardinal_successor_is_cardinal,cardinal,initial_ordinal}.
    $c(\emptyset)\in\cardSucc{c(\emptyset)}$ by
        \obv{cardinal_belongs_to_successor}.
    Hence $\lambda\in\cardSucc{c(\emptyset)}$ by
        \obv{ordinal_transitivity}.
    $r(\emptyset)=\cardSucc{c(\emptyset)}$ by
        \auto{successor_cardinal_map_apply}.
    Hence $\lambda\in r(\emptyset)$.
    Follows by \auto.
\end{proof}
